\documentclass[14pt]{beamer}
\usepackage{omi-template}
\usepackage{hyperref}

\usepackage{calc}
\newlength{\myscalewidth}
\newcommand{\scaletowidth}[2]{%
\setlength{\myscalewidth}{\widthof{#2}}%
\pgfmathsetmacro{\myscaleratio}{#1/\myscalewidth}%
\scalebox{\myscaleratio}{#2}}

\makeatletter
\define@key{beamerframe}{nofills}[true]{% top
  \beamer@frametopskip=0pt\relax%
  \beamer@framebottomskip=0pt\relax%
  \beamer@frametopskipautobreak=\beamer@frametopskip\relax%
  \beamer@framebottomskipautobreak=\beamer@framebottomskip\relax%
  \def\beamer@initfirstlineunskip{%
    \def\beamer@firstlineitemizeunskip{%
      \vskip-\partopsep\vskip-\topsep\vskip-\parskip%
      \global\let\beamer@firstlineitemizeunskip=\relax}%
    \everypar{\global\let\beamer@firstlineitemizeunskip=\relax}}
}
\makeatother


\usepackage{tikz}
\usetikzlibrary{calc}
\usetikzlibrary{shadings}
\usetikzlibrary{arrows, decorations.markings}
\usetikzlibrary{shapes.arrows}
\usetikzlibrary{positioning,fadings,through}

\newcounter{a}
\newcounter{b}

\newcommand{\slice}[5]{
  \pgfmathparse{0.5*#1+0.5*#2}
  \let\midangle\pgfmathresult

  % slice
  \draw[thick,fill=#5] (0,0) -- (#1:1) arc (#1:#2:1) -- cycle;

  % outer label
  %\node[label=\midangle:#4] at (\midangle:0.5) {};

  % inner label
  \pgfmathparse{min((#2-#1-10)/110*(-0.3),0)}
  \let\temp\pgfmathresult
  \pgfmathparse{max(\temp,-0.5) + 0.8}
  \let\innerpos\pgfmathresult
  \node at (\midangle:\innerpos) {\parbox{1in}{\begin{center}#3 \\ \scriptsize #4\end{center}}};
}

\tikzfading[name=fade right,
  left color=transparent!0, right color=transparent!100]
\tikzfading[name=fade left,
  right color=transparent!0, left color=transparent!100]
\tikzfading[name=fade down,
  bottom color=transparent!0, top color=transparent!100]
\tikzfading[name=fade up,
  top color=transparent!0, bottom color=transparent!100]
\tikzfading[name=fade out,
  inner color=transparent!0, outer color=transparent!100]


\def\shadowradius{0.025\textwidth}
\def\zerodistance{1pt}
%
\newcommand\drawshadowbis[1]{
  \begin{pgfonlayer}{shadow}
    \fill[fill=white] ($(#1.south west)+(-\zerodistance,-\zerodistance)$) rectangle ($(#1.north east)+(\zerodistance,\zerodistance)$);

    \begin{scope}
      \clip ($(#1.south west)$) rectangle ++(-\shadowradius,-\shadowradius);
      \fill[fill=white,path fading=fade out]
        ($(#1.south west)$) circle (\shadowradius);
    \end{scope}

    \begin{scope}
      \clip ($(#1.south east)$) rectangle ++(\shadowradius,-\shadowradius);
      \fill[fill=white,path fading=fade out]
        ($(#1.south east)$) circle (\shadowradius);
    \end{scope}

    \begin{scope}
      \clip ($(#1.north west)$) rectangle ++(-\shadowradius,\shadowradius);
      \fill[fill=white,path fading=fade out]
        ($(#1.north west)$) circle (\shadowradius);
    \end{scope}

    \begin{scope}
      \clip ($(#1.north east)$) rectangle ++(\shadowradius,\shadowradius);
      \fill[fill=white,path fading=fade out]
        ($(#1.north east)$) circle (\shadowradius);
    \end{scope}

    \fill[path fading=fade up,fill=white] ($(#1.south west)+((0,-\shadowradius)$) rectangle ($(#1.south east)$);
    \fill[path fading=fade right,fill=white] ($(#1.south east)$) rectangle ($(#1.north east)+((\shadowradius,0)$);
    \fill[path fading=fade down,fill=white] ($(#1.north west)$) rectangle ($(#1.north east)+((0,\shadowradius)$);
    \fill[path fading=fade left,fill=white] ($(#1.south west)$) rectangle ($(#1.north west)+(-\shadowradius,0)$);
  \end{pgfonlayer}
}

\pgfdeclarelayer{shadow}
\pgfsetlayers{shadow,main}

\begin{document}

%%%%%%%%%%%%%%%%%%%%%%%%%%%%%%%%%%%%%%%%%%%%%%%%%%%%%%%%%%%%%%%%
\begin{frame}[nofills]

  \vfill

  \scaletowidth{\textwidth}{\textbf{Ohio's Gateway Mathematics Courses}}

  \begin{center}
    \textit{Bridges to Success Workshop} \\
  April 20 and April 21, 2016
  \end{center}

  \color{dark}

  \vfill
  \hrule
  \vspace{-12pt}

  \begin{center}
    Communication, Outreach and Engagement Co-leads
  \end{center}

  \vspace{-12pt}

  \footnotesize 
  \begin{columns}
    \begin{column}{0.4\textwidth}
      Jim Fowler \\
      The Ohio State University \\
      \texttt{fowler.291@osu.edu}
    \end{column}
    \begin{column}{0.4\textwidth}
      Michelle Younker  \\
      Owens Community College \\
      \texttt{michelle\_younker@owens.edu} \\
    \end{column}
  \end{columns}




\end{frame}

%%%%%%%%%%%%%%%%%%%%%%%%%%%%%%%%%%%%%%%%%%%%%%%%%%%%%%%%%%%%%%%%
\begin{frame}
  \frametitle{Presentation Overview}

  \begin{itemize}
  \item The Ohio Mathematics Initiative
  \item Steering Committee and recommendations
  \item Data
  \item Highlights of work to date
  \item Ohio's gateway mathematics courses
  \end{itemize}
  
\end{frame}

%%%%%%%%%%%%%%%%%%%%%%%%%%%%%%%%%%%%%%%%%%%%%%%%%%%%%%%%%%%%%%%%
\begin{frame}%BADBAD
  \frametitle{Where did the Ohio Mathematics Initiative start?}

  \begin{itemize}
  \item Call for change from institutions \\[0.5ex]
    Concerns about \ldots
    \begin{itemize}
      \item success rates in mathematics courses,
      \item transferability and innovation, \textcolor{dark}{and}
      \item student success.
    \end{itemize}
  \item Statewide Mathematics Summit \hfill\textcolor{dark}{May 2013}
  \item Formation of Steering Committee \hfill\textcolor{dark}{July 2013}
  \end{itemize}
\end{frame}

%%%%%%%%%%%%%%%%%%%%%%%%%%%%%%%%%%%%%%%%%%%%%%%%%%%%%%%%%%%%%%%%
\begin{frame}
  \frametitle{Steering Committee Charge}
  
  To develop expectations and processes \\\quad that result in each campus offering pathways \\\quad that result in mathematics that yield
  \begin{enumerate}
    \item increased success for students \\\quad in the study of mathematics,
    \item a higher percentage of students \\\quad completing degree     
      programs, \textcolor{dark}{and}
    \item effective transferability of credits \\\quad for students moving 
      \\\quad\quad from one institution to another.
  \end{enumerate}
  
\end{frame}

%%%%%%%%%%%%%%%%%%%%%%%%%%%%%%%%%%%%%%%%%%%%%%%%%%%%%%%%%%%%%%%%
\begin{frame}
  \frametitle{Essential Components}

  \begin{enumerate}
    \item Develop high-quality entry-level courses \\ and pathways
    \item Develop transfer policies and processes that foster transfer and encourage innovation
    \item Support constructive engagement of mathematics chairpersons and faculty
    \item Collect, analyze, and share relevant data
    \item Improve student success in college-level mathematics by aligning postsecondary expectations and high school practice
  \end{enumerate}

\parbox{\textwidth}{\footnotesize\color{gray}\href{https://ohiohighered.org/sites/ohiohighered.org/files/uploads/math/MATH-REPORT_FINAL_4.22.14.pdf}{\textit{Rethinking Postsecondary Mathematics.}  Final Report of the Ohio Mathematics Steering Committee, March 2014.}}

\end{frame}

%%%%%%%%%%%%%%%%%%%%%%%%%%%%%%%%%%%%%%%%%%%%%%%%%%%%%%%%%%%%%%%%
\begin{frame}[nofills]

  \vfill  

  \scaletowidth{\textwidth}{``The status quo is unacceptable.''}

  \null\hfill\parbox{0.65\textwidth}{\scriptsize\textcolor{gray}{---\textsc{A Common Vision} for Undergraduate Mathematical Sciences Programs in 2025, Karen Saxe and Linda Braddy, Mathematical Association of America, 2015.}}

  \vfill
\end{frame}

\begin{frame}
  \frametitle{What do we know? \\ Too many entering first-years need remediation.}

\begin{tikzpicture}[overlay,remember picture]
  \node[anchor=north east,rotate=90,yshift=9pt,xshift=-1in] (attribution) at (current page.north east) {\scalebox{0.7}{\scriptsize\textcolor{gray}{Complete College America; used with permission}}};
\end{tikzpicture}


\begin{tikzpicture}
  \node [inner sep=0,text width=0.95\textwidth,rectangle,rounded corners=0pt,fill=white] (box) {\includegraphics[clip,trim=0cm 17.2cm 0cm 4.1cm, width=\textwidth]{too-many-need-remediation.pdf}};
\drawshadowbis{box}
\end{tikzpicture}

  African-American, Hispanic, and low-income students \\
  \quad are more likely to be headed toward remediation.
\end{frame}

\begin{frame}
  \frametitle{What do we know? \\ Very few make it to graduation day}

\begin{tikzpicture}[overlay,remember picture]
  \node[anchor=north east,rotate=90,yshift=9pt,xshift=-1in] (attribution) at (current page.north east) {\scalebox{0.7}{\scriptsize\textcolor{gray}{Complete College America; used with permission}}};
\end{tikzpicture}

\begin{tikzpicture}
  \node [inner sep=0,text width=0.95\textwidth,rectangle,rounded corners=0pt,fill=white] (box) {\includegraphics[clip,trim=0cm 18.6cm 0cm 3cm, width=\textwidth]{very-few-make-it-to-graduation.pdf}};
\drawshadowbis{box}
\end{tikzpicture}

  
\end{frame}


\begin{frame}
  \frametitle{What do we know? \\ Very few make it to graduation day}

\begin{tikzpicture}[overlay,remember picture]
  \node[anchor=north east,rotate=90,yshift=9pt,xshift=-1in] (attribution) at (current page.north east) {\scalebox{0.7}{\scriptsize\textcolor{gray}{Complete College America; used with permission}}};
\end{tikzpicture}

\begin{tikzpicture}
  \node [inner sep=0,text width=0.95\textwidth,rectangle,rounded corners=0pt,fill=white] (box) {\includegraphics[clip,trim=0cm 7.8cm 0cm 13.5cm, width=\textwidth]{very-few-make-it-to-graduation.pdf}};
\drawshadowbis{box}
\end{tikzpicture}

  
\end{frame}

%%%%%%%%%%%%%%%%%%%%%%%%%%%%%%%%%%%%%%%%%%%%%%%%%%%%%%%%%%%%%%%%
\begin{frame}
\frametitle{Students who take College Algebra\ldots}

\begin{columns}
\begin{column}{0.4\textwidth}
\begin{center}
\only<2->{\begin{tikzpicture}[scale=2.7]
    \clip (-1.05,-1.05) rectangle (1.05,1.05);
  \foreach \p/\t/\c in {10/Take Calc I\hspace{1em}\null/blue!25!white, 30/Take business calculus/green!25!white, 60/Do not take any calculus/red!25!white}
  {
    \setcounter{a}{\value{b}}
    \addtocounter{b}{\p}
    \slice{\thea/100*360}
    {\theb/100*360}
    {\p\%}{\t}{\c}
  }
\end{tikzpicture}}\only<1>{%
\begin{tikzpicture}[scale=2.7]
    \clip (-1.05,-1.05) rectangle (1.05,1.05);
  \foreach \p/\t/\c in {10/Take Calc I\hspace{1em}\null/blue!25!white, 90/Do not take Calc I/red!25!white}
  {
    \setcounter{a}{\value{b}}
    \addtocounter{b}{\p}
    \slice{\thea/100*360}
    {\theb/100*360}
    {\p\%}{\t}{\c}
  }
  
\end{tikzpicture}}
\end{center}
\end{column}
\begin{column}{0.4\textwidth}
\uncover<3->{\footnotesize Virtually no students \\\quad who pass college algebra \\\quad ever start Calculus III, \\\quad\quad a key course for STEM.}
\end{column}
\end{columns}

\parbox{\textwidth}{\scriptsize\color{gray}Dunbar, S. 2005. Enrollment flow to and from courses below calculus . In A Fresh State for Collegiate mathematics: Rethinking the Courses below calculus, N.B. Hastings et al. (Eds.). Washington DC: MAA Notes, Mathematical Association of America.}

\end{frame}

%%%%%%%%%%%%%%%%%%%%%%%%%%%%%%%%%%%%%%%%%%%%%%%%%%%%%%%%%%%%%%%%
\begin{frame}[nofills]
  \vfill
  \scaletowidth{\textwidth}{What are we doing?}
  \vfill
\end{frame}

%%%%%%%%%%%%%%%%%%%%%%%%%%%%%%%%%%%%%%%%%%%%%%%%%%%%%%%%%%%%%%%%
\begin{frame}
  \frametitle{Some of the Work to Date \\ \textit{Increasing departmental flexibility regarding pre-requisites
    and credit hours}}

Removed the prescribed pre-requisite requirements \\
\quad from gateway courses for acceptance \\
\quad\quad into the transfer module.

\vfill\hfill\textcolor{dark}{\footnotesize Endorsed late January/early February 2015}
\end{frame}

%%%%%%%%%%%%%%%%%%%%%%%%%%%%%%%%%%%%%%%%%%%%%%%%%%%%%%%%%%%%%%%%
\begin{frame}
  \frametitle{Some of the Work to Date \\
    \textit{Creating a college-level mathematics course definition}}

  A credit-bearing, college-level course in Mathematics \\
  \quad must use the standards required for high school \\
  \quad\quad graduation by the State of Ohio
  as a basis \\
  \quad and must do at least one of the following:
\begin{enumerate}
\item broaden\textcolor{dark}{, or}
\item deepen\textcolor{dark}{, or}
\item extend the student's learning.
\end{enumerate}

\vfill\hfill\textcolor{dark}{\footnotesize Endorsed late January/early February 2014}
\end{frame}

%%%%%%%%%%%%%%%%%%%%%%%%%%%%%%%%%%%%%%%%%%%%%%%%%%%%%%%%%%%%%%%%
\begin{frame}
  \frametitle{Some of the Work to Date \\ \textit{Redesigning Ohio Transfer Module course criteria}}

Revised learning outcomes for\ldots

\begin{itemize}
\item College Algebra \hfill\textcolor{dark}{(\href{http://regents.ohio.gov/transfer/otm/math-stats-log/TMM001-College-Algebra.pdf}{\texttt{TMM001}})}
\item Introductory Statistics\hfill\textcolor{dark}{(\href{http://regents.ohio.gov/transfer/otm/math-stats-log/TMM010.pdf}{\texttt{TMM010}})}
\item Quantitative Literacy\hfill\textcolor{dark}{ (\href{https://www.ohiohighered.org/sites/ohiohighered.org/files/uploads/transfer/documents/OTM/TMM011 Quantitative Reasoning FINALIZED v2- 12-21-2015.pdf}{\texttt{TMM011}})}
\end{itemize}

\vfill\hfill\textcolor{dark}{\footnotesize Endorsed Fall 2015}
  \end{frame}

%%%%%%%%%%%%%%%%%%%%%%%%%%%%%%%%%%%%%%%%%%%%%%%%%%%%%%%%%%%%%%%%
\begin{frame}
  %\frametitle{Examples of Revisions from College Algebra}

  \begin{columns}
    \begin{column}{0.495\textwidth}
      \begin{center}\large\textcolor{gray}{\textbf{Original}}\end{center}

      \textbf{Functions}

      \textbf{1.1} Represent functions verbally, numerically, graphically, and algebraically, including linear, quadratic, polynomial, rational, root/radical/power, exponential, logarithmic and piecewise-defined functions.      
    \end{column}

    \begin{column}{0.01\textwidth}
      \rule{0.01\textwidth}{\textheight}
    \end{column}
    \uncover<2->{
    \begin{column}{0.495\textwidth}
      \begin{center}\large\textcolor{gray}{\textbf{Revised}}\end{center}

      \textbf{1. Functions:} Successful College Algebra students demonstrate a deep understanding of functions whether they are described verbally, numerically, graphically, or algebraically (both explicitly and implicitly).  Students should be proficient working with\ldots
      %the following families of functions:  
      %linear, quadratic, higher-order polynomial, rational, exponential, logarithmic, radical, and piecewise-defined functions (including absolute value).    

    \end{column}
    }
  \end{columns}



\end{frame}

%%%%%%%%%%%%%%%%%%%%%%%%%%%%%%%%%%%%%%%%%%%%%%%%%%%%%%%%%%%%%%%%
\begin{frame}
  \frametitle{\texttt{TMM001} College Algebra}

\footnotesize

\textbf{1. Functions: }
Successful College Algebra students demonstrate a deep understanding of 
functions whether they are described verbally, numerically, graphically, or algebraically (both 
explicitly and implicitly). Students should be proficient working with the following families of functions:  linear, quadratic, higher-order polynomial, rational, exponential, logarithmic, radical, and piecewise-defined functions (including absolute value).

\vfill

The successful College Algebra student can:

\null\hfill\parbox{0.95\textwidth}{\textbf{1a. Analyze functions.}  Routine analysis includes discussion of domain, range, zeros, general function behavior (increasing, decreasing, extrema, etc.). In addition to performing rote processes, the student can articulate reasons for choosing a particular process, recognize function families and anticipate behavior, 
and explain the implementation of a process (e.g., why certain real numbers are excluded from the domain of a given function).}

\end{frame}


%%%%%%%%%%%%%%%%%%%%%%%%%%%%%%%%%%%%%%%%%%%%%%%%%%%%%%%%%%%%%%%%
\begin{frame}
  \frametitle{Some of the Work to Date \\
\textit{Developing mathematics pathways}}

\begin{itemize}
\item Aligned with specific groups of majors
\item Comprised of challenging mathematics content
\item Relevant to the groups of majors
\end{itemize}

\vfill\hfill\textcolor{dark}{\footnotesize Formalized Fall 2015}
\end{frame}

%%%%%%%%%%%%%%%%%%%%%%%%%%%%%%%%%%%%%%%%%%%%%%%%%%%%%%%%%%%%%%%%
\begin{frame}
  \frametitle{What is a math pathway?}

  A \textbf{mathematics pathway} \\
  \quad is a math course \textcolor{gray}{or sequence of courses} taken \\
  \quad by both college-ready and underprepared students \\
  \quad to meet the requirements of their program of study.

  \vfill

  A pathway allows students to \\
  \quad actively engage with mathematical concepts, \\
  \quad access prior knowledge, \textcolor{gray}{and} \\
  \quad reflect on new learning.

  \vfill

  A pathway aligns with specific fields of study.

  \vfill

\scriptsize\vfill\textcolor{dark}{\textit{The New Mathways Project,} Charles A. Dana Center, The University of Texas at Austin, 2015.}
\end{frame}

%%%%%%%%%%%%%%%%%%%%%%%%%%%%%%%%%%%%%%%%%%%%%%%%%%%%%%%%%%%%%%%%
\begin{frame}
  \frametitle{Ohio's Mathematics Pathways}

  Developed in response to the recommendations of the Steering Committee to:
  \begin{itemize}
  \item develop high-quality entry-level courses \\
    \quad and pathways;
  \item increase student success;
  \item make mathematics more relevant to all students;
  \item provide students with \\
    \quad the appropriate mathematics \\
    \quad to successfully support them \\
    \quad in their chosen field of study.
  \end{itemize}
  
\end{frame}

%%%%%%%%%%%%%%%%%%%%%%%%%%%%%%%%%%%%%%%%%%%%%%%%%%%%%%%%%%%%%%%%
\begin{frame}[nofills]
  %\frametitle{Ohio's Mathematics Pathways}

  \begin{tikzpicture}
    \node[anchor=east,circle,fill=red!50!white,draw=red!60!white] (1) {1};
    \node [anchor=west,fill=red!50!white, single arrow, draw=red!60!white, single arrow head indent=0ex, single arrow head extend=6pt] at (0,0) {Statistics Pathway\hspace{1cm}\null};
  \end{tikzpicture}

  College-level introductory statistics courses \\
  \quad designed for students \\
  \quad\quad without a Calculus background \\
  \quad and who \\
  \quad\quad do not require College Algebra or Calculus.

  \vfill
  \hrule
  \vfill

  Part of the general education requirement \\
  \quad for majors in fields that may include the following: \\[-2ex]
  \begin{columns}
    \color{dark}
    \begin{column}{0.35\textwidth}
      \begin{center}
        Nursing \\ Nutrition
        \end{center}
    \end{column}
    \begin{column}{0.65\textwidth}
      \begin{center}
        Social Work \\ Associates in Business
        \end{center}
    \end{column}
  \end{columns}

  \vfill

\end{frame}

%%%%%%%%%%%%%%%%%%%%%%%%%%%%%%%%%%%%%%%%%%%%%%%%%%%%%%%%%%%%%%%%
\begin{frame}[nofills]
  \begin{tikzpicture}
    \node[anchor=east,circle,fill=blue!30!white,draw=blue!40!white] (2) {2};
    \node[anchor=west,fill=blue!30!white, single arrow, draw=blue!40!white, single arrow head indent=0ex, single arrow head extend=6pt] at (0,0) {Quantitative Reasoning Pathway\hspace{1cm}\null};
  \end{tikzpicture}

  College-level courses designed to emphasize \\
  \quad quantitative thinking and problem solving \\
  \quad using quantitative methods.

  \vfill
  \hrule
  \vfill

  Part of the general education requirement \\
  \quad for majors in fields that may include the following: \\[-2ex]
  \begin{columns}
    \color{dark}
    \begin{column}{0.5\textwidth}
      \begin{center}
        Communication \\
        Criminal Justice \\
        Fine arts
        \end{center}
    \end{column}
    \begin{column}{0.65\textwidth}
      \begin{center}
        Education \footnotesize\\
        (Elementary, History, \\
        Social Studies, etc.)
        \end{center}
    \end{column}
  \end{columns}

  \vfill

\end{frame}

%%%%%%%%%%%%%%%%%%%%%%%%%%%%%%%%%%%%%%%%%%%%%%%%%%%%%%%%%%%%%%%%
\begin{frame}[nofills]
  \begin{tikzpicture}
    \node[anchor=east,circle,fill=green!50!white,draw=green!60!white] (3) {3};
    \node[anchor=west,fill=green!50!white, single arrow, draw=green!60!white, single arrow head indent=0ex, single arrow head extend=6pt] at (0,0) {STEM Preparation Pathway\hspace{1cm}\null};
  \end{tikzpicture}

  College-level courses designed \\
  \quad for students in mathematics-intensive majors.

  \vfill

  \textcolor{dark}{Examples: College Algebra, Pre-Calculus, \\
    \quad Trigonometry, Business Calculus, and Calculus.}

  \vfill
  \hrule
  \vfill

  Part of the general education requirement \\
  \quad for majors in fields that may include the following: \\[-2ex]
  \begin{columns}
    \color{dark}
    \begin{column}{0.5\textwidth}
      \begin{center}
            Business			 \\
            Engineering \\
            Physics

        \end{center}
    \end{column}
    \begin{column}{0.65\textwidth}
      \begin{center}
        Chemistry \\[0.5ex]
        Education \footnotesize\\
        (Math, Science, \\ Technology, etc.)
        \end{center}
    \end{column}
  \end{columns}

  \vfill

\end{frame}

%%%%%%%%%%%%%%%%%%%%%%%%%%%%%%%%%%%%%%%%%%%%%%%%%%%%%%%%%%%%%%%%
\begin{frame}[nofills]
  \vfill
  \scaletowidth{\textwidth}{Thank You!}
  \vfill
\end{frame}

%%%%%%%%%%%%%%%%%%%%%%%%%%%%%%%%%%%%%%%%%%%%%%%%%%%%%%%%%%%%%%%%
\begin{frame}
  \frametitle{Contact Information \\ Communication, Outreach, \& Engagement Subgroup}

  Jim Fowler, Co-lead \\
  The Ohio State University \\
  \texttt{fowler.291@osu.edu} \\
  773.809.5659 

  \vfill
  
    Michelle Younker, Co-lead  \\
     Owens Community College \\
     \texttt{michelle\_younker@owens.edu} \\
     567.661.2640

     \vfill
\end{frame}

%%%%%%%%%%%%%%%%%%%%%%%%%%%%%%%%%%%%%%%%%%%%%%%%%%%%%%%%%%%%%%%%
\begin{frame}
  \frametitle{Contact Information \\ Ohio Department of Higher Education}


Dr. Paula Compton \\
    Associate Vice Chancellor, Articulation and Transfer \\
    \texttt{pcompton@highered.ohio.gov}

    \vfill
    
Hideo Tsuchida \\
    Director of Articulation and Transfer Policy \\
    \texttt{htsuchida@highered.ohio.gov}

    \vfill
    
\end{frame}

%%%%%%%%%%%%%%%%%%%%%%%%%%%%%%%%%%%%%%%%%%%%%%%%%%%%%%%%%%%%%%%%
\begin{frame}[allowframebreaks]
  \frametitle{Resources and References}

  \begin{itemize}
\item \href{https://ohiohighered.org/sites/ohiohighered.org/files/uploads/math/MATH-REPORT_FINAL_4.22.14.pdf}{Rethinking Postsecondary Mathematics: Final Report of the Ohio Mathematics Steering Committee}
\item \href{https://ohiohighered.org/mathematics-initiative-documents}{OMI Website}
\item \href{https://www.ohiohighered.org/transfer/transfermodule/learningoutcomes}{OTM Guidelines and Learning Outcomes}
\item \href{https://www.ohiohighered.org/content/ohio_mathematics_initiative_resources}{OMI Newsletter, Fact Sheets, and Reading List}
\item \href{https://www.ohiohighered.org/mathematics-initiative-resources/presenter-request}{OMI Speaker Request Form}
\item \href{http://www.completecollege.org/docs/Ohio_remediation.pdf}{Complete College America Remediation}

  \end{itemize}


\end{frame}


\end{document}

